\bookchapter{Capitolo 13: App di utilità consigliate}

\chapterintro{Questo capitolo raccoglie strumenti di produttività, salute e viaggio, oltre a funzioni di sistema e multi-dispositivo: pulizia e anti-spam come Qingzhuo, blocchi integrati del telefono, meteo e allerta terremoti; controllo remoto del PC con ToDesk/Sunlogin, ufficio AI con Trae, sfondi animati, trasferimento rapido in LAN con LocalSend, appunti multi-dispositivo, codici dinamici con Aegis/Microsoft Authenticator, agenda con TickTick e scansione di documenti cartacei.}

\booksection{Strumenti di produttività: gestisci tempo e attenzione}

Gli strumenti di produttività sono lo «scheletro» della vita digitale: ti aiutano a mettere ordine tra idee, attività e materiali sparsi. Qui sotto le raccomandazioni per scenario.

\booksubsection{Liste di cose da fare (gestione attività)}

\begin{itemize}
  \item TickTick Per chi: ha bisogno di registrare in modo unificato attività di lavoro, vita e studio insieme agli impegni; uno dei consigliati principali di questa sezione. Punti di forza: sincronizzazione tra telefono/computer/tablet/Web; liste + calendario + pomodoro in un'unica app; l'input in linguaggio naturale come “riunione domani alle 15” o “consegna del report ogni lunedì” riconosce automaticamente data e ora e le regole di ripetizione; supporta sotto-attività, etichette, priorità, allegati ed elenchi condivisi. Perfetto per una gestione integrata di “agenda + spunta delle attività + ripasso quotidiano”.
  \item Come usarlo meglio TickTick: Casella di raccolta: butta tutto in «Oggi/Inbox» e la sera classifica nelle liste lavoro, famiglia, studio.
  \item Vista calendario: trascina riunioni e scadenze sul calendario per vedere i conflitti, molto più intuitivo della sola lista.
  \item Attività ricorrenti: medicine, bollette, palestra impostale come regole, senza reinserirle ogni volta.
  \item Timer pomodoro: avvia la concentrazione direttamente sull'attività e a fine settimana sai davvero quanto hai completato.
  \item Integrazione con il telefono: all'assistente vocale puoi dire «Ricordami…» per creare promemoria di sistema; per le scadenze importanti attiva anche il calendario di sistema a doppia garanzia.
\end{itemize}

Todoist Per chi: amanti del minimalismo e utenti multi-piattaforma. Punti di forza: interfaccia sobria, gerarchia chiara (attività sotto-attività sotto-sotto-attività), punti karma per costruire abitudini.

Microsoft To Do Per chi: ufficio e utenti Outlook. Punti di forza: completamente gratuito, senza pubblicità; «La mia giornata» si svuota ogni giorno e alleggerisce la mente; integrazione profonda con l'ecosistema Microsoft.

Note/Promemoria di sistema Per chi: utenti leggeri che non vogliono app di terze parti. Punti di forza: esperienza nativa, zero ritardo, aggiunta vocale con Siri/Xiao AI, integrazione con i widget di sistema.

\booksubsection{Timer pomodoro (allenamento della concentrazione)}

\begin{itemize}
  \item Forest Per chi: chi non riesce a stare lontano dal telefono. Punti di forza: durante la concentrazione pianti un albero; se esci dall'app presto, l'albero muore; usa la «vita virtuale» per creare feedback positivi e puoi piantare insieme agli amici.
  \item Tide Per chi: ha bisogno del rumore bianco per concentrarsi o addormentarsi. Punti di forza: suoni naturali (pioggia, onde del mare, foresta) + pomodoro + meditazione in un'unica app, con un'interfaccia minimale e rilassante.
  \item Principio: classici 25 minuti di focus + 5 minuti di pausa; puoi adattare a 50/10 o 90/20 secondo il tuo ritmo.
\end{itemize}

\booksubsection{Calendari (visualizzare il tempo)}

\begin{itemize}
  \item Calendario di sistema Per chi: tutti. Punti di forza: notifiche native più stabili, consumo di batteria minimo, aggiunta vocale con Siri/Xiao AI.
  \item Google Calendar Per chi: utenti multi-dispositivo e chi condivide l'agenda con altri. Punti di forza: la funzione «Obiettivi» ti pianifica automaticamente palestra/studio; integrazione intelligente con Gmail per estrarre eventi dalle email.
\end{itemize}

\booksubsection{App di note (consolidare la conoscenza)}

\begin{itemize}
  \item Notion Per chi: chi vuole gestire insieme note, attività, database e wiki; studenti, freelance e team. Punti di forza: knowledge base completa, editor a blocchi, puoi creare homepage personali, liste di lettura, board di progetto; ricco ecosistema di template.
  \item Obsidian Per chi: gestori avanzati di conoscenza, ricercatori, scrittori. Punti di forza: archiviazione Markdown locale, link bidirezionali + grafo delle relazioni, le tue note diventano un «secondo cervello»; i dati restano completamente tuoi.
  \item flomo Per chi: chi vuole catturare idee frammentarie senza scrivere saggi. Punti di forza: note a scheda, come pubblicare un post; supporta etichette \# e la navigazione cronologica.
  \item Note Apple Per chi: utenti iPhone/Mac con esigenze di registrazione leggere. Punti di forza: avvio rapido, scansione documenti/disegni a mano/elenchi, sincronizzazione iCloud, gratuito e profondamente integrato con il sistema.
  \item Evernote Per chi: utenti esperti con molte ritaglio web e archiviazione. Punti di forza: potente estensione di web clipper, OCR del testo nelle immagini, storico e stabile; la versione premium è a pagamento.
\end{itemize}

\booksubsection{Strumenti PDF}

\begin{itemize}
  \item PDF Expert Per chi: utenti iPad/iPhone, professionisti che annotano contratti, studenti. Punti di forza: considerato il miglior editor PDF su iOS, annotazioni fluide, firma, unione/divisione, compilazione moduli.
  \item WPS Per chi: chi lavora con Word/Excel/PPT/PDF. Punti di forza: suite da ufficio completa, ottima compatibilità con cloud cinesi e file WeChat; la versione gratuita basta per l'uso quotidiano.
  \item Adobe Acrobat Per chi: chi ha requisiti elevati sullo standard PDF (legale, design, editoria). Punti di forza: standard ufficiale, conversione in PDF/A impeccabile, sincronizzazione multi-dispositivo e firma elettronica professionali.
\end{itemize}

\booksubsection{Scansione (digitalizzare documenti cartacei)}

Contratti cartacei, ricevute, compiti, appunti su lavagna, fotocopie di documenti d'identità… la fotocamera del telefono sostituisce già la maggior parte degli scanner. Flusso: scatta chiaramente ritaglio e raddrizzamento automatici esporta PDF/immagine OCR opzionale.

\begin{itemize}
  \item CamScanner Per chi: professionisti e studenti che scansionano spesso contratti, scontrini e documenti. Punti di forza: riconoscimento automatico dei bordi, correzione prospettica, OCR, PDF multipagina e sincronizzazione cloud.
  \item Microsoft Lens Per chi: studenti e chi deve trasformare lavagne/pagine di libro in Word. Punti di forza: gratuito e senza pubblicità, esportazione in un clic su OneNote/Word, riconoscimento testo molto accurato.
  \item Scansione in Note Apple Per chi: utenti iPhone con esigenze di scansione leggere. Punti di forza: senza installare app, basta toccare l'icona fotocamera in Note; rimozione automatica delle ombre, firma e annotazioni; gratuito e comodo.
  \item Baimiao / CamScanner Web / «Scansione documenti» della fotocamera di marca: Huawei, Xiaomi, OPPO, vivo hanno quasi tutte una modalità «Documento/Scansione» integrata; su Android puoi usare Google Camera o la scansione di sistema. Utile per scattare al volo una pagina di contratto.
  \item Consigli di scatto: luce uniforme, telefono parallelo al foglio, punta ai quattro angoli; i documenti multipagina scattali in ordine e poi unisci in PDF; per i documenti sensibili attenti a filigrana e posizione di salvataggio, e dopo l'uso eliminali definitivamente da «Eliminati di recente».
\end{itemize}

\booksection{Salute e sport}

Le app di salute sono il «assistente personale» più sottovalutato del telefono: tracciano continuamente sonno, frequenza cardiaca e attività, e raccontano lo stato di lungo periodo meglio di una visita saltuaria.

\booksubsection{Fitness}

\begin{itemize}
  \item Keep Per chi: allenamento a casa, principianti, chi vuole dimagrire e tonificare. Punti di forza: molti corsi gratuiti, dall'HIIT allo yoga, piani generati con AI, community vivace.
  \item Nike Training Club (NTC) Per chi: chi cerca un coach professionale e un allenamento sistematico. Punti di forza: prodotto ufficiale Nike, corsi gratuiti, dimostrazioni di coach professionisti, classificati per gruppo muscolare/attrezzatura.
  \item Fitness Apple Per chi: utenti Apple Watch e chi vuole chiudere gli anelli di attività. Punti di forza: tre anelli giornalieri (movimento/allenamento/degenza), obiettivi visivi, integrazione con i dati della frequenza cardiaca dell'orologio.
\end{itemize}

\booksubsection{Monitoraggio del sonno}

\begin{itemize}
  \item AutoSleep Per chi: utenti Apple Watch che vogliono conoscere la proporzione sonno profondo/leggero. Punti di forza: tracciamento automatico dormendo con l'orologio, senza avvio/arresto manuale, punteggio di qualità del sonno e consigli per andare a letto.
  \item Sleep Cycle Per chi: chi non dorme con l'orologio e vuole usare il telefono. Punti di forza: analizza colpi di tosse e russamento con il microfono, sveglia intelligente nella fase di sonno leggero per un risveglio meno traumatico.
  \item Xiaomi Fitness Per chi: utenti di braccialetti/orologi Xiaomi. Punti di forza: gratuito e integrato con i dispositivi Xiaomi, tracciamento di frequenza cardiaca/saturazione/sonno, ottimo rapporto qualità-prezzo.
\end{itemize}

\booksubsection{Meditazione e rilassamento}

\begin{itemize}
  \item Calm Per chi: professionisti stressati e chi ha difficoltà ad addormentarsi. Punti di forza: meditazione guidata «Daily Calm» + storie per dormire (lette da celebrità), interfaccia elegante, ricchi contenuti in inglese.
  \item Tide Per chi: utenti cinesi che cercano il rumore bianco per concentrarsi o dormire. Punti di forza: app cinese di alta qualità, suoni naturali + corsi di meditazione + pomodoro, con la guida in cinese ancora più familiare.
  \item Xiao Sleep Per chi: ha bisogno di suoni che aiutino a dormire e vuole registrare i discorsi nel sonno o i russamenti. Punti di forza: una quantità enorme di rumori bianchi da sovrapporre a piacere (pioggia + ventola + falò), spegnimento a timer prima di dormire, e la registrazione dei discorsi nel sonno è divertente e utile.
\end{itemize}

\booksubsection{Salute femminile}

\begin{itemize}
  \item Flo Per chi: chi registra il ciclo, cerca una gravidanza o segue la salute femminile. Punti di forza: previsioni AI accurate di ciclo/ovulazione, community anonima, design curato.
  \item Meiyou Per chi: utenti donne in Cina che vogliono una community e contenuti divulgativi in cinese. Punti di forza: contenuti localizzati molto ricchi, con conoscenze dedicate a ogni fase: ciclo, ricerca gravidanza, gravidanza e cura del bambino.
\end{itemize}

\booksubsection{Promemoria farmaci}

\begin{itemize}
  \item Assistente farmaci Per chi: pazienti cronici che assumono farmaci a lungo termine e figli che devono gestire le medicine degli anziani. Punti di forza: riconosce i farmaci scansionando il codice a barre, crea automaticamente la lista dei farmaci, invia promemoria a orario e registra le assunzioni, che puoi condividere con i familiari.
\end{itemize}

\booksubsection{Hub dei dati}

\begin{itemize}
  \item Salute Apple / Google Fit Per chi: tutti gli utenti smartphone. Punti di forza: raccoglie tutti i dati da telefono, braccialetto e app di terze parti (passi, frequenza cardiaca, sonno, ciclo), è il «centro unico della cartella sanitaria».
\end{itemize}

\booksection{Viaggi e spostamenti}

In viaggio, poche app giuste ti fanno evitare innumerevoli grattacapi. Qui sotto le raccomandazioni per l'intero flusso di viaggio.

\booksubsection{Voli e treni ad alta velocità}

\begin{itemize}
  \item Umetrip Per chi: professionisti che volano spesso. Punti di forza: dati ufficiali dell'aviazione civile, la fonte più autorevole per stato dei voli/cambi di gate/allerte di ritardo; mostra anche il modello dell'aereo e la mappa dei posti.
  \item 12306 Per chi: tutti i viaggiatori in treno in Cina. Punti di forza: unico canale ufficiale di acquisto, supporta liste d'attesa (compra automaticamente se c'è un rimborso), gestione comoda di biglietti studenteschi/di gruppo.
  \item VariFlight Per chi: chi è sensibile alla puntualità dei voli e deve accompagnare o prelevare qualcuno all'aeroporto. Punti di forza: previsioni dei ritardi molto accurate; puoi condividere il viaggio con la famiglia, che vede in tempo reale la tua posizione e l'ora di arrivo.
\end{itemize}

\booksubsection{Mappe e navigazione}

\begin{itemize}
  \item Amap (Gaode Maps) Per chi: chi vive nelle città cinesi, automobilisti e chi prende i taxi. Punti di forza: traffico in tempo reale molto accurato ed efficace nell'evitare la congestione; taxi e condivisione della corsa in un'unica app; ricchissime le opzioni di trasporto pubblico con cambi.
  \item Baidu Maps Per chi: chi ha bisogno di informazioni sui luoghi e navigazione indoor. Punti di forza: street view / navigazione AR, mappe dettagliate di centri commerciali e attrazioni, ricchi pacchetti vocali.
  \item Google Maps Per chi: chi viaggia all'estero. Punti di forza: copertura globale completa, recensioni locali autentiche, street view e mappe offline sono un salvaventura all'estero.
\end{itemize}

\booksubsection{Traduzione}

\begin{itemize}
  \item Google Traduttore Per chi: viaggiatori e studenti di lingue. Punti di forza: traduzione bidirezionale in tempo reale in «Conversazione», traduzione istantanea con la «Fotocamera»; indispensabile in viaggio, supporta pacchetti offline.
\end{itemize}

\booksubsection{Conversione valutaria}

\begin{itemize}
  \item Minimal Exchange Rates Per chi: chi viaggia, acquista da siti esteri o fa shopping cross-border. Punti di forza: interfaccia essenziale, oltre 160 valute, tassi aggiornati in tempo reale, confronto simultaneo di più valute, conversione rapida premendo a lungo sul numero.
\end{itemize}

\booksubsection{Meteo}

\begin{itemize}
  \item Meteo di sistema / Weather.com.cn Per chi: tutti, prima scelta per previsioni e allerte ufficiali. Punti di forza: le app meteo dei produttori e i canali ufficiali come «Meteo Cina» si collegano all'osservatorio; allerte per pioggia intensa, tifoni, caldo estremo vengono inviate con banner/notifiche, alta autorevolezza.
  \item Caiyun Weather Per chi: ha bisogno di previsioni di pioggia precise al minuto. Punti di forza: l'accuratezza su “pioverà nelle prossime due ore?” è tra le migliori del settore, la mappa dinamica della fascia di pioggia è intuitiva e arriva la notifica push prima che inizi a piovere.
  \item Carrot Weather Per chi: appassionati di meteo che cercano ironia e widget ricchi. Punti di forza: meteo letto da un AI sarcastica, moltissimi widget personalizzabili, mappe radar professionali; la versione a pagamento è più completa.
  \item Punti chiave sulle allerte: le app meteo devono avere le notifiche attive per ricevere le allerte; le allerte sono blu, gialla, arancione, rossa; il rosso indica il rischio meteorologico più alto e conviene adattare i programmi. Per altre modalità di attivazione vedi la sezione 7 di questo capitolo «Allerte meteo / Allerte calamità naturali / Allerta terremoti».
\end{itemize}

\booksubsection{Guide di viaggio}

\begin{itemize}
  \item Mafengwo Per chi: chi ama guide dettagliate, viaggi in auto e itinerari profondi. Punti di forza: ricchi racconti originali degli utenti, guide su cosa fare/percorsi/errori da evitare, scaricabili offline.
  \item Xiaohongshu Per chi: giovani viaggiatori che cercano esperienze autentiche e locali di nicchia/cibo. Punti di forza: formato visivo con foto e video brevi, contenuti come «3 giorni a XX città» aggiornati e puntuali, ottima ricerca.
\end{itemize}

\booksection{Introduzione agli strumenti di automazione}

L'automazione è una leva: «una configurazione, risparmio a lungo termine»; può azzerare le decine di azioni ripetitive di ogni giorno. Questa sezione è solo un'introduzione ai concetti; per usi concreti e tutorial dettagliati vedi Capitolo 23: Pensiero computazionale e introduzione all'automazione.

\begin{itemize}
  \item «Scorciatoie» iOS: lo strumento di automazione ufficiale di Apple, senza programmazione; con un tap puoi eseguire azioni combinate come «a casa: attiva Wi-Fi + abbassa luminosità + leggi il meteo».
  \item Automazione Android: le scene di sistema (come «Automazioni» di Xiaomi, «Smart Scene» di Huawei) sono le più semplici da usare; Tasker è potente ma ha una curva di apprendimento ripida, adatto agli smanettoni.
  \item Idea di base: identifica prima le «azioni che ripeti ogni giorno», poi chiediti «posso sostituirle con una condizione di attivazione + più azioni?». Esempi dettagliati e usi avanzati nel Capitolo 23.
\end{itemize}

\booksection{Browser e motori di ricerca: il tuo ingresso alle informazioni}

Il browser è la tua finestra su internet, il motore di ricerca la bussola nell'oceano delle informazioni. Usarli bene ti fa risparmiare molti errori.

\booksubsection{Come funziona un browser:}

Quando digiti un indirizzo nel browser succede questo: (1) il browser interroga il server DNS per l'IP del dominio; (2) invia una richiesta HTTP/HTTPS a quel server; (3) il server restituisce il codice della pagina (HTML, CSS, JavaScript); (4) il browser interpreta il codice e lo rende nella pagina che vedi.

\booksubsection{Browser consigliati:}

\begin{itemize}
  \item Microsoft Edge: basato sul motore Chromium, veloce e compatibile. Include l'assistente AI Copilot per riassumere pagine, tradurre e conversare. Supporta la sincronizzazione multi-dispositivo; consigliatissimo come browser principale.
  \item Safari: browser predefinito su Apple, profondamente integrato con il sistema, efficiente e risparmia batteria. Ottima privacy, ma poche estensioni.
  \item Chrome: il più usato al mondo, ricco di estensioni. Però consuma molta memoria e sul telefono consuma più batteria.
\end{itemize}

\booksubsection{Scegliere il motore di ricerca:}

\begin{itemize}
  \item Bing: motore di ricerca Microsoft, integrato con Edge. I riassunti AI danno direttamente le risposte e non solo i link: ricerca molto efficiente. Buona anche la ricerca per immagini.
  \item Yandex: il più grande motore russo, molto forte nella ricerca immagini e nell'immagine inversa. Utile per trovare immagini simili e riconoscere oggetti.
  \item Baidu: principale motore per le informazioni in Cina, ma con molti annunci; serve capacità di discernimento.
  \item Google: il motore più preciso al mondo, ma in Cina richiede una rete speciale.
\end{itemize}

\booksubsection{Guida alla ricerca efficiente di informazioni:}

\begin{itemize}
  \item Usa bene la ricerca AI: i riassunti AI di Bing e Copilot di Edge rispondono direttamente alle domande e ti aiutano a estrarre le informazioni chiave, molto più efficiente che aprire ogni link.
  \item Riconoscere false informazioni e annunci: nei risultati di ricerca le voci marcate «Annuncio» o «Sponsorizzato» sono promozioni a pagamento, non risultati organici. Valuta l'autorevolezza della fonte e preferisci siti ufficiali e media noti.
  \item Scaricare le app correttamente: digita «nome app + sito ufficiale» nel motore di ricerca, ad esempio «WeChat sito ufficiale», e preferisci i risultati con dominio ufficiale. Non installare software da store di terzi: spesso includono software indesiderato.
  \item Come riconoscere i siti ufficiali: per scaricare Windows 11 visita microsoft.com, per Chrome visita google.com. Controllare che il dominio corrisponda al marchio è il metodo più basilare.
\end{itemize}

\booksubsection{Conoscere i domini dei siti web:}

Il dominio è il «numero civico» del sito: leggerlo ti aiuta a capirne natura e provenienza:

\begin{itemize}
  \item Significato dei suffissi più comuni: .com — organizzazione commerciale, il più comune
  \item .cn — dominio nazionale della Cina
  \item .org — organizzazione non profit
  \item .net — servizi di rete
  \item .edu — istituzione educativa
  \item .gov — istituzione governativa
  \item .xyz/.tech/.info — domini generici di nuova generazione, bassa soglia di accesso, usare con prudenza
\end{itemize}

Domini nazionali/territoriali: .jp — Giappone

.kr — Corea

.uk — Regno Unito

.de — Germania

.tw — Taiwan

.hk — Hong Kong

Suggerimento di sicurezza: i siti con suffisso .gov e .edu sono di solito più affidabili perché la registrazione è più rigorosa. I domini economici come .xyz e .top sono spesso usati da siti di phishing: stai attento.

\booksubsection{Blocco annunci:}

\begin{itemize}
  \item AdGuard Per chi: tutti gli utenti infastiditi da popup, annunci e pubblicità nelle app. Punti di forza: blocco annunci a livello di sistema (non solo browser), filtra banner e tracciamento nelle app, risparmia dati e velocizza le pagine; la versione a pagamento copre tutti i dispositivi.
\end{itemize}

\booksubsection{Browser per la privacy:}

\begin{itemize}
  \item Brave Per chi: chi tiene alla privacy e non vuole essere tracciato dalla pubblicità. Punti di forza: basato su Chromium, blocca annunci e tracker fin da subito, include la modalità Tor, più veloce di Chrome e con meno consumo di memoria.
  \item DuckDuckGo Per chi: chi non vuole che una ricerca scateni pubblicità ovunque. Punti di forza: motore di ricerca omonimo (non registra la cronologia), cancellazione di tutte le tracce con un tap («Fire Button»), esperienza coerente su tutte le piattaforme.
\end{itemize}

\booksubsection{Confronto tra browser:}

\begin{center}
\begin{tabularx}{\linewidth}{|X|X|X|}
\hline
\textbf{Browser} & \textbf{Punti di forza} & \textbf{Per chi} \tabularnewline
\hline
Microsoft Edge & Motore Chromium + Copilot AI integrato: riassume pagine, traduce, conversa & Professionisti orientati all'efficienza, utenti intensivi di AI (consigliato come principale) \tabularnewline
Safari & Integrazione profonda con Apple, efficiente, risparmia batteria, ottima privacy & Utenti iPhone/Mac che tengono a autonomia e integrazione \tabularnewline
Chrome & Il più usato al mondo, ricchissimo di estensioni & Utenti avanzati che usano molte estensioni e sincronizzazione multi-dispositivo \tabularnewline
Brave & Blocca annunci e tracker fin da subito, più veloce & Chi tiene alla privacy e vuole bloccare gli annunci senza installare estensioni \tabularnewline
DuckDuckGo & Motore di ricerca omonimo senza cronologia, cancellazione tracce con un tap & Utenti sensibili alla privacy che non vogliono tracciamento personalizzato \tabularnewline
\hline
\end{tabularx}
\end{center}

\begin{tcbpractice}{Prova a creare un'automazione «Rilassarsi prima di dormire»}
\begin{enumerate}
  \item Crea la condizione di attivazione: un orario specifico (es. 22:30).
  \item Aggiungi le azioni: abbassa volume > attiva Non disturbare > riproduci playlist rilassante > abbassa luminosità.
  \item Salva e prova.
\end{enumerate}

\end{tcbpractice}

\booksection{Strumenti di sistema utili: pulizia, anti-spam e sicurezza}

Oltre alle app di produttività e salute, ciò che influisce davvero sulla vita quotidiana sono tre cose: «meno disturbi, spazio sufficiente, truffatori fuori». Questa sezione elenca prima i software utili più comuni, poi approfondisce le funzioni di blocco integrate nel telefono: molte persone usano lo stesso telefono da anni senza aver mai attivato queste opzioni.

\booksubsection{Pulizia di sistema e ottimizzazione dello spazio}

\begin{itemize}
  \item Qingzhuo Per chi: utenti Android che vogliono pulire con leggerezza cache, file residui e cartelle vuote senza farsi disturbare dalla “suite tutto in uno” dei gestori di telefono. Punti di forza: leggera e senza pubblicità, supporta la scansione approfondita dei residui delle app, l'individuazione dei file di grandi dimensioni e l'analisi dello spazio occupato; più misurata dei commerciali “accelerazione con un tocco”, ideale per utenti esperti che fanno manutenzione periodica. Attenzione: prima di pulire verifica i percorsi; non selezionare alla cieca le cartelle di WeChat, del rullino fotografico e dei documenti; per le cache di sistema privilegia lo strumento di pulizia integrato del produttore e usa gli strumenti di terze parti solo come complemento.
  \item Pulizia dello spazio integrata Per chi: tutti, soprattutto anziani e chi non vuole complicazioni. Punti di forza: Android [Impostazioni] > [Archiviazione] > [Pulisci e velocizza]; iOS [Impostazioni] > [Generali] > [Spazio su iPhone]. Il sistema conosce meglio la struttura della cache delle proprie app e il rischio di cancellare per errore è minimo: deve essere la prima scelta.
  \item Gestione spazio integrata in WeChat/QQ Per chi: chi usa molto la chat. Punti di forza: WeChat [Io] > [Impostazioni] > [Generali] > [Spazio di archiviazione], pulizia mirata per chat e tipo di file; più sicura degli strumenti esterni «pulisci WeChat».
  \item Principio: se esiste la funzione di sistema o dell'app, non usare terzi; se puoi selezionare per percorso/tipo, non fare «pulisci tutto con un tap». L'obiettivo è liberare spazio, non cancellare dati importanti.
\end{itemize}

\booksubsection{Software anti-spam e di sicurezza}

\begin{itemize}
  \item National Anti-Fraud Center Per chi: tutti, soprattutto anziani, chi lavora in finanza e chi fa spesso acquisti online o riceve chiamate sconosciute. Punti di forza: prodotto ufficiale, avvisi chiamate in arrivo, rischio SMS, controllo URL sospetti; integrato con il database anti-frode della polizia, l'anello più autorevole delle «tre linee di difesa».
  \item Centro sicurezza del marchio / Phone Manager Per chi: utenti del marchio corrispondente. Punti di forza: Huawei Phone Manager, Xiaomi Security Center, OPPO Phone Manager, vivo Phone Manager ecc., integrano blocco chiamate, antivirus, gestione permessi e protezione pagamenti; di solito gratuiti e meglio integrati col sistema. Sconsigliato installare una seconda suite commerciale per evitare conflitti di permessi.
  \item Servizi anti-spam dell'operatore Per chi: utenti bombardati da chiamate automatizzate e vocali AI. Punti di forza: China Mobile invia KTFSR al 10086; China Telecom segui «Tianyi Anti-Spam»; China Unicom imposta in «Unicom Mobile Manager». Il blocco avviene lato rete e ferma molti numeri senza far squillare il telefono.
  \item AdGuard Per chi: tutti gli utenti infastiditi da popup, annunci e pubblicità nelle app. Punti di forza: blocco annunci a livello di sistema (non solo browser), filtra banner e tracciamento nelle app, risparmia dati e velocizza le pagine; la versione a pagamento copre tutti i dispositivi.
  \item Attenzione alle suite «pulizia/blocco» di terzi: alcuni restano in background e chiedono continuamente permessi su rubrica e SMS; la loro capacità di blocco non è necessariamente superiore e aumentano consumo e rischi privacy. Preferisci le funzioni integrate del sistema e poi l'app anti-frode ufficiale.
\end{itemize}

\textbf{Principi di scelta del software utile}\par

\begin{itemize}
  \item Prima il sistema: per blocco, pulizia e permessi attiva prima le funzioni integrate.
  \item Integrazione ufficiale: National Anti-Fraud Center e anti-spam dell'operatore sono canali ufficiali e si possono sovrapporre.
  \item Terzi con parsimonia: per la stessa funzione basta un'app, per evitare conflitti di permessi e consumo di risorse.
  \item Meno è meglio: più strumenti installi, più notifiche e permessi aumentano e più è probabile essere disturbati dagli strumenti stessi.
\end{itemize}

\booksubsection{Blocco chiamate indesiderate integrato: come attivarlo per marca}

Qui sotto i percorsi per marca. I nomi possono variare leggermente con la versione del sistema: puoi cercare direttamente nelle impostazioni «indesiderate» «blocco» «chiamate sconosciute». Complementare al Capitolo 8 «Gestione delle chiamate indesiderate»: qui l'accento è sull'attivazione, lì sulla segnalazione dei numeri e sui ricorsi.

\begin{center}
\begin{tabularx}{\linewidth}{|X|X|X|}
\hline
\textbf{Piattaforma} & \textbf{Percorso} & \textbf{Opzioni consigliate} \tabularnewline
\hline
iPhone / iOS & {}[Impostazioni] > [Telefono] / [Messaggi] & Silenzia chiamate sconosciute; contatti bloccati; in Messaggi «Filtra mittenti sconosciuti» \tabularnewline
Xiaomi / Redmi & App Telefono > Impostazioni > Blocco chiamate indesiderate; oppure Centro sicurezza & Blocco intelligente, squillo unico per numeri sconosciuti, filtro parole chiave SMS \tabularnewline
Huawei / Honor & Telefono > Blocco chiamate indesiderate; Phone Manager > Blocco chiamate indesiderate & Riconoscimento chiamate, lista nera, blocco intelligente SMS \tabularnewline
OPPO / realme / OnePlus & App Telefono > Blocco chiamate indesiderate; oppure Telefono > Blocco chiamate indesiderate & Blocco chiamate, riconoscimento truffe, regole di blocco SMS \tabularnewline
vivo / iQOO & Telefono > Blocco chiamate indesiderate; i Manager & Blocco chiamate indesiderate, avvisi truffe, filtro messaggi \tabularnewline
Samsung & Telefono > Impostazioni > Blocca numeri / Identificazione chiamate e spam & Blocca numeri sconosciuti, riconoscimento spam (supporto diverso per area) \tabularnewline
\hline
\end{tabularx}
\end{center}

\begin{tcbinfo}{Consigli per combinare tre livelli di protezione}
Uso quotidiano consigliato: «blocco di sistema + anti-spam operatore + National Anti-Fraud Center». Il sistema gestisce marcature e liste nere, l'operatore blocca le chiamate automatizzate lato rete, l'anti-frode aggiunge gli avvisi ufficiali. Non serve impostare tutto al massimo, altrimenti blocchi anche corrieri, food delivery e medici: puoi lasciare passare i contatti e i contatti recenti.

\end{tcbinfo}

\begin{tcbpractice}{Esercizio pratico: attiva oggi stesso il blocco chiamate indesiderate}
\begin{enumerate}
  \item Apri le impostazioni dell'app Telefono, trova «Blocco chiamate indesiderate / Silenzia chiamate sconosciute» e attiva il blocco intelligente.
  \item Controlla anche il telefono dei tuoi genitori: stesso percorso + attiva l'anti-spam dell'operatore (China Mobile: invia KTFSR al 10086).
  \item Installa «National Anti-Fraud Center» dallo store, registrati e attiva gli avvisi chiamate.
  \item Prova a chiamare un numero di test già segnalato (o controlla il registro dei blocchi) per verificare che funzioni.
\end{enumerate}

\end{tcbpractice}

\booksection{Allerte meteo, calamità naturali e terremoti}

Il blocco chiamate ferma le «truffe telefoniche», le allerte fermano «calamità e maltempo estremo». Molti installano solo l'app meteo per vedere la temperatura, senza notifiche né allerta terremoti: nei secondi prima di un'allerta rossa per pioggia intensa, un tifone o un terremoto, il telefono può guadagnarti tempo prezioso per metterti al sicuro.

\booksubsection{Allerte meteo: come ricevere pioggia intensa, tifoni e caldo estremo}

\begin{itemize}
  \item Prima i canali ufficiali: le allerte dell'Amministrazione Meteorologica Cinese / Osservatorio Meteorologico Centrale / osservatori provinciali sono le più autorevoli. Puoi seguire l'account WeChat ufficiale «Meteo Cina», Weibo, o scaricare app ufficiali come «Meteo Cina».
  \item App Meteo di sistema: Huawei Weather, Xiaomi Weather, OPPO Weather, vivo Weather, «Meteo» di Apple ecc., la maggior parte invia le allerte dell'osservatorio locale; nelle impostazioni dell'app Meteo assicurati che le notifiche «Allerte meteorologiche / Maltempo» siano attive.
  \item Caiyun Weather e simili per la pioggia al minuto: utili per decidere prima di uscire o andare a prendere i bambini se serve l'ombrello; per le allerte resta decisivo l'osservatorio ufficiale.
  \item Significato dei colori delle allerte (generale): Blu: stai attento, di solito puoi uscire normalmente e seguire gli aggiornamenti.
  \item Giallo: rafforza le precauzioni, riduci le attività all'aperto non necessarie.
  \item Arancione: grave, resta in casa se possibile, tieniti lontano dalle zone a rischio.
  \item Rosso: particolarmente grave, segui le indicazioni di sospensione scuole/lavoro e di evacuazione; non forzare strade allagate, colate detritiche o vento forte.
\end{itemize}

Tipi di allerta comuni: pioggia intensa, nevicata intensa, tifone, caldo estremo, ondata di freddo, vento forte, fulmini, grandine, ghiaccio sulle strade, nebbia fitta, siccità, rischio di frane.

\booksubsection{Allerte per calamità naturali: non solo «pioggia»}

\begin{itemize}
  \item Chi le emette: gli osservatori per il meteo; frane e colate sono spesso emesse congiuntamente da protezione civile, gestione acque e dipartimenti di emergenza; anche i governi locali pubblicano tramite radio di emergenza, SMS, 12379 e altri canali.
  \item Come riceverle da più canali: Notifiche del sistema + allerte app Meteo.
  \item Segui su WeChat gli account ufficiali locali di «Protezione Civile» e «Ufficio Meteorologico».
  \item Canali 12379 per le allerte nazionali (SMS/servizi, a seconda dell'attivazione locale).
  \item TV, radio di emergenza, altoparlanti del villaggio/quartiere — in montagna o durante la stagione delle piogge non contare solo sul telefono.
\end{itemize}

Avvisi per attività all'aperto: durante un tifone non avvicinarti a cartelloni e non passare sotto gli alberi; con pioggia intensa non entrare in garage interrati, gallerie o sottopassi; in montagna con allerta «rischio geologico» allontanati subito da pendii e valli.

Integrazione con la guida di primo soccorso: in caso di smarrimento dopo un disastro, telefono scarico o necessità di localizzazione, vedi Capitolo 32 Scenari terremoto/calamità e Capitolo 7 Comunicazioni satellitari e senza rete.

\booksubsection{Allerta terremoti: non è la stessa cosa della «previsione»}

L'allerta terremoti (EEW) sfrutta il fatto che le onde elettromagnetiche viaggiano più velocemente delle onde sismiche e che le stazioni captano prima le onde P per inviare avvisi alle aree colpite secondi o decine di secondi prima dell'arrivo delle onde più distruttive (onde S). Non prevede se ci sarà un terremoto, ma può dare il tempo di «mettersi al sicuro, spegnere il fuoco, uscire dall'ascensore». Più sei vicino all'epicentro, meno tempo hai, e a volte non basta.

\begin{itemize}
  \item Soluzioni ufficiali e di settore: in Cina continentale i principali servizi includono la Rete Nazionale di Allerta Terremoti e i servizi promossi da istituti come il Chengdu Institute of Disaster Prevention; alcuni produttori preinstallano o possono ricevere allerte online.
  \item App di terze parti: cerca nello store app ufficiali/istituzionali come «Allerta terremoti» (prima di installare controlla lo sviluppatore ed evita le app contraffatte che chiedono permessi inutili).
  \item Cosa fare dopo l'allerta: in casa riparati vicino (sotto il tavolo, angoli di muri portanti, mobili solidi), proteggi testa e collo; tieniti lontano da finestre, pareti esterne e mobili alti; non buttarti sulle scale o dalla finestra; se sei in ascensore scendi al piano più vicino; all'aperto vai in spazio aperto, lontano da cavi, cartelloni e edifici pericolosi. Dopo il terremoto controlla gas e impianto elettrico e attenzione alle repliche.
  \item Non silenziare: l'allerta terremoti deve poter «bucare» silenzio e Non disturbare; di notte tieni il telefono a portata di mano e non mettere il volume a 0.
\end{itemize}

\booksubsection{Allerta terremoti integrata: come attivarla per marca}

I percorsi qui sotto possono variare con la versione del sistema: cerca direttamente nelle impostazioni «terremoto» «allerta» «allerta sicurezza». iOS ha un supporto limitato per le allerte terremoti in Cina: di solito serve un'app di terze parti affidabile e permessi di localizzazione e notifiche attivi.

\begin{center}
\begin{tabularx}{\linewidth}{|X|X|X|}
\hline
\textbf{Piattaforma} & \textbf{Percorsi comuni} & \textbf{Note} \tabularnewline
\hline
Xiaomi / Redmi & Phone Manager / Centro sicurezza > Allerta terremoti; oppure cerca «terremoto» nelle impostazioni & Alcune versioni hanno il servizio integrato, serve connessione e localizzazione \tabularnewline
Huawei / Honor & Impostazioni > Sicurezza > Notifiche di emergenza / Allerta terremoti & Il nome può essere «Notifiche di emergenza», include anche i terremoti \tabularnewline
OPPO / realme / OnePlus & Impostazioni > Sicurezza (SOS emergenza) > Allerta terremoti & Alcuni modelli dipendono dalla copertura del servizio locale \tabularnewline
vivo / iQOO & App Meteo oppure Impostazioni > Sicurezza/Altro > Allerta terremoti & Fai riferimento al menu effettivo della versione di sistema \tabularnewline
iPhone / iOS & Di solito non c'è un interruttore completo integrato; installa un'app di allerta terremoti affidabile e consenti le notifiche & L'app deve poter girare in background e avere la localizzazione; Non disturbare non deve silenziarne gli avvisi \tabularnewline
Samsung e altri & Impostazioni > Sicurezza e situazioni di emergenza / Avvisi di emergenza wireless & La copertura degli «avvisi di emergenza» varia per area e operatore \tabularnewline
\hline
\end{tabularx}
\end{center}

\begin{tcbwarn}{Falsi miti sull'allerta terremoti}
\begin{itemize}
  \item Allerta $\neq$ previsione: non ti dice «ci sarà un terremoto domani», avvisa solo quando il terremoto è già avvenuto e le onde non sono ancora arrivate.
  \item Non ovunque c'è molto preavviso: vicino all'epicentro può esserci solo qualche secondo o non esserci affatto; l'allerta serve soprattutto alle aree circostanti.
  \item Notifiche spente = app inutile: anche la migliore allerta non arriva se silenzi, uccidi il processo o neghi la localizzazione.
  \item App contraffatte: diffida di software di «previsione terremoti» che chiedono rubrica e SMS; preferisci le soluzioni di sistema o istituzionali.
\end{itemize}

\end{tcbwarn}

\begin{tcbinfo}{Come combinare le capacità di allerta}
Consiglio quotidiano: «notifiche allerte meteo di sistema + canali ufficiali meteo/protezione civile + allerta terremoti (di sistema o app affidabile)». Come per l'anti-spam, non serve installare decine di app simili; l'importante è notifiche attive, localizzazione concessa e Non disturbare che non blocchi gli avvisi. Dopo aver attivato il blocco chiamate sul telefono dei genitori, controlla anche allerte meteo e terremoti.

\end{tcbinfo}

\begin{tcbpractice}{Esercizio pratico: 10 minuti per controllare le impostazioni di allerta}
\begin{enumerate}
  \item Apri l'app Meteo di sistema, vai nelle impostazioni e attiva le notifiche «Allerte meteorologiche / Maltempo».
  \item Cerca nelle impostazioni di sistema «terremoto» o «allerta», trova e attiva allerta terremoti/emergenza (Android); su iPhone installa un'app di allerta terremoti affidabile e consenti le notifiche.
  \item Segui gli account WeChat dell'ufficio meteorologico o della protezione civile locale e consenti i push.
  \item Verifica che il suono dell'allerta non venga silenziato dal Non disturbare; controlla anche il telefono dei tuoi genitori.
\end{enumerate}

\end{tcbpractice}

\booksection{Collaborazione multi-dispositivo: controllo remoto, trasferimento rapido e sicurezza account}

Un telefono funziona bene soprattutto quando «gioca in squadra» con PC e tablet. Questa sezione aggiunge strumenti di uso frequente: controllo remoto del PC dal telefono, ufficio AI e collaborazione da remoto, sfondi animati, trasferimento rapido in LAN, appunti multi-dispositivo, codici dinamici 2FA e integrazione tra scansione cartacea e agenda. Per un approccio più sistematico al multi-dispositivo vedi Capitolo 17.

\booksubsection{Controllo remoto del PC dal telefono: Sunlogin / ToDesk}

\begin{itemize}
  \item Sunlogin Remote Control Per chi: chi deve usare da remoto un Windows o Mac di casa/ufficio; spesso usato anche per assistere da remoto i genitori (vedi Capitolo 27). Punti di forza: ecosistema locale maturo, l'app telefono controlla il desktop, trasferisce file, accende/spegne (in parte serve hardware); la versione gratuita basta per un uso personale leggero.
  \item ToDesk Per chi: chi cerca connessione stabile e bassa latenza per lavoro/assistenza da remoto. Punti di forza: connessione con codice dispositivo + password, telefono controlla PC e PC controlla telefono (a seconda della versione), trasferimento file e gestione multi-monitor fluidi; i nodi cinesi sono di solito buoni.
  \item Uso di base: installa il client sul PC e accedi/resta online installa l'app sul telefono avvia la sessione con ID dispositivo + password di connessione (o lo stesso account) usa touchpad/gesti del telefono per controllare il desktop.
  \item Limiti di sicurezza (da ricordare): Se uno sconosciuto si spaccia per «assistenza/forze dell'ordine/operatore» e ti chiede di installare Sunlogin o ToDesk per controllare da remoto il tuo telefono o computer, rifiuta sempre: è una truffa tipica (vedi Capitolo 27).
  \item Per uso personale: attiva l'autenticazione a due fattori, cambia la password predefinita, esci dall'account o disattiva l'accesso non presidiato quando non usi; sui computer aziendali rispetta le policy di sicurezza.
  \item In luoghi pubblici attento a qualcuno che sbircia schermo e password.
\end{itemize}

\booksubsection{Software di ufficio AI: Trae / WorkBuddy / Dumate}

\begin{itemize}
  \item Trae Per chi: sviluppatori e collaborazione su documenti, ambiente amichevole in cinese. Punti di forza: IDE/assistente di programmazione AI di ByteDance, la modalità SOLO genera un progetto con una frase; efficiente per scrivere codice e revisionare documenti. Se gira sul PC, dal mobile puoi seguire/controllare i compiti con il controllo remoto o le funzioni ufficiali (fai riferimento all'ultima versione del prodotto). Vedi Capitolo 41 Panorama degli strumenti di programmazione AI.
  \item WorkBuddy Per chi: professionisti che vogliono l'AI per verbali di riunione, email, report e documenti d'ufficio. Punti di forza: assistente AI per scenari da ufficio, adatto a «bozza—revisione—riepilogo»; alcuni flussi si avviano dal telefono e il lavoro pesante si fa sul PC.
  \item Dumate Per chi: team e individui che hanno bisogno di collaborazione e gestione integrata di documenti e conoscenza. Punti di forza: strumento di ufficio AI collaborativo, con raccolta materiali e editing condiviso; dal telefono si viste e modifiche leggere, per la formattazione complessa meglio il grande schermo del PC.
  \item Cosa significa «controllare il PC dal telefono»: molte funzioni AI girano su PC/web; dal telefono, collegandoti al PC con Sunlogin/ToDesk o usando l'app/web del prodotto, puoi «pilotare» i compiti AI sul PC anche in metro. Principio: calcolo pesante e documenti lunghi sul PC; controllare lo stato e modificare due righe dal telefono.
  \item Privacy: prima di caricare contratti, documenti d'identità o dati dei clienti leggi le condizioni privacy del prodotto e le policy aziendali; per i materiali aziendali usa gli strumenti designati dall'azienda.
\end{itemize}

\booksubsection{Sfondi animati: Live Desktop di miHoYo (telefono + PC)}

\begin{itemize}
  \item Rengong Zhuomian Per chi: amante dei personaggi miHoYo che cerca sfondi animati gratuiti. Punti di forza: app di sfondi animati gratuita lanciata da miHoYo, disponibile per telefono e PC; i personaggi in Live2D/animati possono diventare il tuo desktop e, su alcuni modelli/sistemi, anche lo sfondo della schermata di blocco o del desktop.
  \item Batteria e prestazioni: gli sfondi animati consumano più batteria e GPU delle immagini statiche; con batteria scarica o modelli datati usali solo in carica o abbassa il frame rate (se l'app lo permette).
  \item Origine dell'installazione: preferisci sito ufficiale/store ufficiali, evita le «versioni crackate» di terzi che nascondono trojan; per la versione PC attento al percorso di installazione e all'avvio automatico.
\end{itemize}

\booksubsection{Trasferimento rapido in LAN: LocalSend / Feige / AirDrop}

\begin{itemize}
  \item LocalSend Per chi: chi mescola Android/Windows/iOS/Mac/Linux e vuole trasferimenti cross-platform senza cloud pubblico. Punti di forza: open source, connessione diretta in LAN; foto, video e documenti sullo stesso Wi-Fi, senza cavi e senza passare dai server dei produttori; la velocità dipende dal router.
  \item Feige Chuanshu Per chi: uffici con PC Windows collegati tra loro e vecchi ambienti di lavoro su rete locale. Punti di forza: storico strumento di trasferimento e chat su LAN, ideale per inviare file di grandi dimensioni in intranet; scaricalo solo da canali ufficiali e sulla rete aziendale rispetta le policy IT.
  \item AirDrop Per chi: utenti iPhone, iPad e Mac. Punti di forza: condivisione istantanea nell'ecosistema Apple, scoperta Bluetooth + trasferimento Wi-Fi diretto, senza consumare dati; puoi impostare «Solo contatti» per ridurre i disturbi.
  \item Altre soluzioni: Mi Share di Xiaomi, Huawei Share, vivo Share ecc. sono rapidi tra dispositivi dello stesso marchio; «File Transfer Helper» di WeChat è comodo ma limitato per dimensione e dispositivi di accesso; per file grandi preferisci LocalSend/AirDrop.
  \item Suggerimenti: entrambi sullo stesso Wi-Fi; su Wi-Fi pubblici trasferire file sensibili è rischioso, preferisci hotspot privati; al termine verifica sul dispositivo ricevente che il file sia integro.
\end{itemize}

\booksubsection{Sincronizzazione appunti multi-dispositivo: KDE Connect / collegamento telefono-PC}

\begin{itemize}
  \item «Appunti universali» Apple Per chi: utenti iPhone + Mac. Punti di forza: con lo stesso Apple ID e Handoff attivo, copi sul telefono incolli sul Mac (e viceversa), per testo e alcuni file.
  \item KDE Connect Per chi: Android + Windows/Linux/Mac, chi vuole sincronizzazione appunti open source, mirror notifiche e input da remoto. Punti di forza: dopo l'associazione in LAN sincronizza gli appunti, invia le notifiche del telefono al PC, usa il telefono come telecomando/slideshow/touchpad; i dati non passano da account di terzi.
  \item «Collegamento telefono-PC» dei marchi: Huawei Multi-Screen Collaboration/Super Device, Xiaomi Mi Share, vivo Share, OPPO Connect, Honor Connect ecc., spesso con appunti multi-schermo e trascinamento file; l'esperienza è migliore con lo stesso marchio.
  \item Canale ufficiale Windows + Android: «Collegamento al telefono / Link to Windows» (evoluzione di Your Phone) permette di vedere notifiche e foto del telefono sul PC e, in alcuni scenari, copiare e incollare; su Samsung e altri è più integrato.
  \item Sicurezza: negli appunti possono restare password e codici di verifica; su computer condivisi evita la sincronizzazione globale, dopo l'uso non reincollare nelle caselle password; anche alcuni metodi di input leggono gli appunti: attento ai permessi.
\end{itemize}

\booksubsection{Codici dinamici 2FA: Aegis / Microsoft Authenticator}

\begin{itemize}
  \item Perché usare un autenticatore invece dell'SMS: i codici SMS possono essere intercettati con lo swap SIM; i codici TOTP sono generati in locale con un algoritmo temporale, non dipendono dal segnale e sono più resistenti a phishing e intercettazione (restano da proteggere da trojan e screenshot).
  \item Aegis Authenticator Per chi: utenti Android che preferiscono open source e controllo locale. Punti di forza: open source e gratuito, backup crittografato e importazione/esportazione, interfaccia pulita, sblocco dell'elenco token con impronta/faccia.
  \item Microsoft Authenticator Per chi: utenti account Microsoft, Office 365/aziendali; su iOS e Android. Punti di forza: approvazione login con push, monitoraggio password e backup account (a seconda della versione), integrazione con Windows e l'identità aziendale.
  \item Google Authenticator / 2FAS e simili: comuni nell'ecosistema; l'importante è avere un backup crittografato per non perdere i token cambiando telefono.
  \item Cosa fare quando cambi telefono: migra dal vecchio dispositivo o usa i «codici di backup/chiavi di recupero» del servizio; stampa i codici di backup su carta e conservali in due posti, non solo sullo stesso telefono.
  \item Con i gestori password: alcuni gestori password memorizzano anche i TOTP; è comodo ma concentra il rischio; per gli account di alto valore consigliamo «gestore password + autenticatore separato». Vedi Capitolo 14.
\end{itemize}

\booksubsection{Combinazioni consigliate: come organizzare la giornata}

\textbf{Esempio di flusso di lavoro}\par

\begin{itemize}
  \item Mattina: controlla l'agenda di oggi in TickTick dai un'occhiata a meteo/allerte apri Trae/documenti sul PC e metti il telefono da parte.
  \item Giornata: i documenti cartacei trasformati in PDF con app di scansione, inviati al PC con LocalSend/AirDrop; i testi lunghi bozzettati sul PC con strumenti di ufficio AI.
  \item Fuori casa: collegati al PC con ToDesk/Sunlogin per continuare le attività; per l'accesso ad account importanti usa i codici dinamici di Aegis/Microsoft Authenticator.
  \item Sera: segna come completate le voci in TickTick; se vuoi rilassarti cambia sfondo con Live Desktop.
\end{itemize}

\begin{tcbwarn}{Regole d'oro per controllo remoto e sicurezza account}
\begin{itemize}
  \item Se qualcuno ti chiede di installare software di controllo remoto riaggancia e contatta tu l'assistenza ufficiale.
  \item Non dire a nessuno codici di verifica o codici dinamici, nemmeno a «corrieri» o «assistenza clienti».
  \item Prima di usare controllo remoto o installare software AI su computer aziendali, chiedi al reparto IT se è consentito.
  \item Su reti pubbliche sii particolarmente cauto con trasferimenti file e accessi da remoto.
\end{itemize}

\end{tcbwarn}

\begin{tcbpractice}{Esercizio pratico: fai funzionare una volta il «multi-dispositivo»}
\begin{enumerate}
  \item Collega telefono e PC allo stesso Wi-Fi e trasferisci una foto con LocalSend (o AirDrop).
  \item Crea in TickTick un'attività «Stasera ore 20: organizzare le scansioni» e verifica che sia visibile sia su telefono che su PC.
  \item Attiva il 2FA su email o WeChat (se supportato), collegalo con Aegis o Microsoft Authenticator e salva i codici di backup.
  \item Installa Sunlogin o ToDesk sul tuo computer, connettiti una volta con il tuo telefono, poi esci o rafforza la password: impara il processo senza lasciare buchi di sicurezza.
\end{enumerate}

\end{tcbpractice}
